\section{Introduction}
\label{sec:intro}

A backdoored classifier behaves normally on clean inputs and predicts an
attacker's class whenever a trigger appears in the input
\cite{li2022backdoorsurvey}. The attacker plants the behavior by poisoning a
small share of the training data (often 1\% or less) and the defender receives
the trained model without ever seeing that data. Input-level detection asks the
most practical question in this setting. Given the deployed model and one
incoming input, does that input carry a trigger? A detector that answers well
needs no retraining and runs in front of the deployed model.

Prediction Shift Backdoor Detection (PSBD) \cite{li2025psbd} was proposed to find
poisoned samples in a training set. Its score needs only the trained model and a
small clean set, so we apply it to each incoming input instead.
It rests on the observation that a backdoored input is decided by a shortcut the
model learned quickly and firmly, while a clean input is decided by a broad set
of features. When the network is perturbed at inference, the broad decision
wobbles and the shortcut holds. PSBD therefore runs the model a few times with
dropout active, measures how much the probability of the predicted class falls
and flags the inputs whose prediction barely moves. It needs only a small clean
validation set to set its threshold, and it reports strong results on ResNet-18
and VGG-16.

Nothing in that argument depends on convolutions, yet the perturbation does not
carry over unchanged. The original work applies dropout after the residual
addition of each ResNet block, before the activation. A transformer block has
no such default and no activation after its residual additions. The input to
attention, the input to the MLP, each branch output, the stream after each
addition, the attention heads and the MLP hidden units are all different tensors
with different roles, and dropout is not the only thing one can inject there,
since whole tokens or channels can also be masked.

Adapting PSBD to a Vision Transformer (ViT) \cite{dosovitskiy2021an} is
therefore a design problem with two parts, where to perturb and what to inject.
We show that this choice decides whether the detector works. Masking whole
tokens at the input of every attention block, which we call PSBD-TM, reaches a
mean AUROC of \HeadlineAurocAdaptive{}, against \PublishedAurocAdaptive{} for
dropout on the residual stream, our adaptation of the original site, which we
call PSBD-RD. We then explain why. A backdoor in a ViT is a single direction in
the residual stream, written late in the network and carried by attention from
the trigger's own tokens to the class token, so removing those tokens before
attention removes the route.

\paragraph{Contributions.} We make three claims.

\begin{enumerate}
  \item A placement that makes PSBD work on transformers. Over \PanelCellsCached{}
    backdoored ViT-B/16 models on \PanelDatasetsCached{} datasets, PSBD-TM beats
    PSBD-RD by \HeadlineGainAdaptiveAuroc{} [\HeadlineGainAdaptiveAurocLow{},
    \HeadlineGainAdaptiveAurocHigh{}] AUROC, by \GainsLowestRate{} at
    \PanelRateLow{} poisoning and by \SwinGainRecommendedMinusPublished{} over
    \SwinCells{} Swin-S models (\cref{sec:results,sec:swin}).
  \item A comparison with \DetectorsComparedCount{} input-level detectors ported to
    the same models, splits and false-positive budgets, in which PSBD-TM ranks
    1st of \DetectorsAurocDefensesRanked{} defenses and PSBD-RD ranks
    5th (\cref{sec:robustness:detectors}).
  \item A mechanistic account, tested on the same models, of why the attention
    input wins, why whole tokens beat coordinates and why noise at that site is
    weaker than masking (\cref{sec:mechanism}).
\end{enumerate}

Every number in this paper is produced by a script from the saved results, and
the build fails on a number typed by hand. The code and the generators are
released. \Cref{sec:limitations} states what the evidence does not cover.
